\begin{table*}[!ht]
\centering
\caption{Analysis of Related Works and Primary Studies Cited}
\label{tab:related_works}
\footnotesize
% As larguras somam ~0.91\textwidth, deixando espaço para o padding (tabcolsep)
\begin{tabular}{p{0.15\textwidth} p{0.38\textwidth} p{0.38\textwidth}}
\toprule
\textbf{Ref.} & \textbf{Research Gaps and Limitations found} & \textbf{Future Works} \\ \midrule

\cite{usman_2022} &
Absence of formal reference architecture, operational complexity of fragmentation, telemetry limitations, and scalability and correlation challenges. & 
Intelligent Automation via AI/ML, Security integration, Enhancing Developer Experience, and Unified Observability Platforms. \\ \addlinespace

\cite{araujo_2024} &
Lack of holistic reference architectures, resource constraints (overhead) in fog nodes, and high heterogeneity of edge devices. & 
Lightweight AI/ML for local anomaly detection (TinyML), standardized interoperability protocols, and adaptive observability frameworks. \\ \addlinespace

\cite{zanatta_23} &
Fragmented focus on metrics over logs/traces in 5G research, high memory/CPU overhead of standard CNOT stacks, and lack of standardized log structures. & 
Integration of distributed tracing, research on geographically distributed B5G setups, scalability analysis of COTS ML, and observability in Open RAN components. \\ \addlinespace

\cite{kosinska_2023} &
Lack of unified terminology, insufficient practical automation via AI/ML (AIOps gap), and poor adaptation of tools for edge computing environments. & 
Development of observability maturity models, deeper integration with eBPF for low-overhead visibility, and alignment with FinOps and GreenOps. \\ \addlinespace

\cite{faseeha_2025} &
Conceptual confusion between monitoring and observability pillars, lack of practical multi-pillar approaches, and difficulties in tracing failures in highly dynamic environments. & 
Evolution of frameworks for modern deployment paradigms (service mesh, event-driven), automation of data analysis, and optimization of observability overhead. \\ \addlinespace

\cite{han_2024} &
Inconsistent problem definitions for RCA and high cloud-native heterogeneity across traces, metrics, and logs. & 
Improving multi-modal data correlation at scale and enhancing Graph Attention Networks for dynamic topologies. \\ \addlinespace

\cite{gomes_2025} &
Difficulty in correlating multi-modal data (metrics, logs, traces), inadequacy of traditional monitoring for sophisticated CNAs, and complex root cause analysis (RCA) in distributed environments. & 
Advanced ML-based resource management, automated diagnostic mechanisms integrated with auto-scaling, and research on low-overhead real-time data collection. \\ \addlinespace

Wei et al. \cite{wei_2025} &
Lack of a systematic bridge between desirable system qualities and the concrete cloud-native practices and tools used to achieve them, and the initial validation scope is limited (expert-based, not covering all contexts). &
Provide mechanisms/interfaces to extend and continuously evolve the map (keeping practices/tools up to date) and perform broader validations beyond the initial expert review. \\ \addlinespace

\cite{deng_2024} &
Missing clear research roadmap for cloud-native lifecycles, high cognitive load for developers (DevEx), and insufficient real-time security observability. & 
Integration of SecOps with performance telemetry, improvement of Developer Experience (DevEx), and observability optimization for serverless and green computing. \\ \bottomrule
\end{tabular}
  \par\medskip\ABNTEXfontereduzida\selectfont\textbf{Fonte:} Written by author (2026) \par\medskip
\end{table*}